\chapter{The Emergence of Behavior}
% OUTLINE Ch3 -- THE BUILT EXHIBIT, the one chapter of this face that pays by
% THEOREM, not shape. Exhibit = the period-3 residue-cycle (CycleOfThree.lean,
% verified claim-by-claim by both chairs, source-READ): cycle_of_three (line 40,
% by decide) proves order EXACTLY three (closes-at-3 AND period!=1 AND period!=2);
% residue_invariant_under_cycle (line 50, cases<;>rfl) fixes the residue; and the
% 3 is EARNED NON-CIRCULARLY as |pair (+) post| = 2(+)1 (coproduct cardinality via
% Sum, proved bijection to the three-fold) -- the preface's FENCEPOST earned.
% SIX RIDERS (both chairs, held): (1) AXIOM-CLAIM OFF THE BUILD -- proofs are by
% decision-procedure + structural identity, the machine DISPLAYS its axiom basis;
% no bare "axiom-free" boast beyond #print axioms. (2) STAND ON THE THEOREMS, not
% the #eval display strings. (3) NO Lean identifiers in prose. (4) electron/null/
% positron = the MODEL'S three-fold states, illustration-not-identity + electron-
% scope, not a real-particle claim. (5) LAB-GAP where the residue reads a number.
% (6) "emergence" model-internal (defined by its own theorem), Anderson left out.
% Unabashed math + riders; bold-math/fenced-referent; step-not-move; bracket-not-
% point. Gate: Kodo (earn/citation, exhibit pre-checked) + Beastmaster (arrival-
% read: STANDS-ON-THE-THEOREM held hardest).

The last two chapters read shapes. This one reads a theorem, and the difference
is the whole reason the chapter exists. The strain rings and the taper damps it,
and under the damping most of what the forced point stirred up decays away to the
floor. The question this chapter answers is what does \emph{not} decay --- what
survives the coarsening of the first chapter and the damping of the second and
comes back, unchanged, every time. That surviving thing is the emergent behavior
of the model's fluid face, and unlike the strain and the taper it stands on a
result the machine has proved.

Begin with the structure the surviving behavior lives in. The model's argument
carries, at its core, a three-fold: three states the fine account cycles
through, which the earlier volumes met under several names --- the calibration's
rotating roles, the three colors of a closed loop, and, in the reading nearest
the model's subject, the triple electron / null / positron. A single operation,
call it the \emph{fold}, carries each state to the next around the ring; write
$\rho$ for it, and $\rho^{k}$ for the fold applied $k$ times. To read the
emergent behavior is to ask what the fold does under repetition: whether the
states wander, or settle, or return.

They return, and they return on a schedule the machine has settled exactly. The
fold brings every state back to itself after \emph{three} applications and after
no fewer:
\[
  \rho^{3} = \mathrm{id}, \qquad \rho^{1} \neq \mathrm{id}, \qquad
  \rho^{2} \neq \mathrm{id}.
\]
All three of these the machine proves together, not by inspection and not by
assertion but by its own decision procedure: it checks, exhaustively over the
finite three-fold, that three folds restore each state and that one and two folds
do not, and it settles the matter with a yes. The order of the cycle is three,
proved to be three and proved to be neither one nor two. And a companion result,
proved this time by plain structural identity --- the two sides reduce to the
same expression --- carries the reading the chapter is after: the residue the
whole apparatus reports is \emph{invariant} under the full turn, unchanged by the
three folds that leave every state where it began. That invariance is the
behavior that does not decay. It is not damped away with the strain, because it
is the one thing the cycle returns intact; it emerges from the fold the way a
standing pattern emerges from a rotation that closes.

\begin{fence}
The machine does not merely run these folds and report what it saw; it
\emph{proves} the return, and it is made to display, alongside the proof, the
axioms the proof rests on. The reading of this chapter leans on that displayed
basis and on the decision procedure and the structural identity that do the work
--- never on a bare claim of purity beyond what the machine prints. What is
claimed is exactly what is proved: order three, not one and not two, and a
residue invariant under the turn.
\end{fence}

There is one more thing the machine proves, and it is what makes the three
honest rather than circular. It would be no achievement to announce that a
three-state cycle has period three; the three would have been smuggled in by
counting the states. The apparatus does not do that. It earns the three from
below, out of two separate and simpler things: a \emph{pair} --- the plus and
minus of the model, two states and no more --- and a single \emph{closing post},
the null that shuts the ring. The count of the whole is the count of the pair
joined to the count of the post,
\[
  \lvert\, \text{pair} \sqcup \text{post} \,\rvert
    \;=\; 2 \;\oplus\; 1 \;=\; 3,
\]
where the join is the coproduct of the two --- a disjoint union of separate
things --- and the addition is not a postulated operation but the size of that
union, counted off the constructed whole. And the machine proves that the
three-fold \emph{is} that union: it exhibits a correspondence between them and
proves it a bijection in both directions, so that the three-fold's period is the
union's size and not its own constructor-count. The three is therefore the
preface's fencepost, arrived at honestly: two --- the pair --- plus the one
closing post, earned as a cardinality and never read off the answer it was meant
to explain.

That is the emergent behavior of this face, and it is worth being exact about the
word, because this book uses it in one sense and guards the others out. Emergence
here is not a gesture at the mystery of the many-from-the-few, and it borrows no
authority from the physics of complex systems. It is defined entirely by the
result just read: the emergent behavior is the structure that survives the
coarsening and the damping and returns unchanged under the cycle --- the
invariant of a fold proved to close at three. Nothing more is meant by it, and
nothing about real collective phenomena is claimed through it. What emerges is
what the theorem says survives, and the theorem is the whole of the claim.

Two fences close the chapter. The three states --- the electron, the null, and
the positron the fold cycles through --- are the \emph{model's} own, three states
of the model's account wearing those names; nothing here measures, models, or
claims a real electron or a real positron, and the reading does not widen past
the one model the series has spoken of throughout. And the residue this cycle
holds invariant is read, where it is read as a number at all, to the model's own
resolution --- the instrument's own value, never the laboratory's, here as in
every chapter. The behavior that emerges has been read, and it has been read off
a theorem: the next chapter turns from what the model builds to the wall it does
not cross.
